\section{Introduction}\label{sec:introduction}

An integer Lipschitz function on the triangular lattice is a discrete
surface with a hard constraint on its nearest-neighbor slopes.  We weight
each unit height change by a parameter $x$.  The associated contours form
the loop $O(2)$ model on the hexagonal lattice.  We prove that the smoothly
averaged surface has a Gaussian free field limit throughout
$1/\sqrt2\le x\le1$, including the predicted critical endpoint.  The
normalization is independent of the domain and of its lattice approximation.

\subsection{Model, approximation, and theorem}

Let $\hexlat$ be the hexagonal lattice embedded in $\R^2$ so that its face
centers form the triangular lattice
\[
 \T=\Z(1,0)+\Z(1/2,\sqrt3/2).
\]
Two faces are adjacent when they share an edge.  We identify each face with
its center whenever using the triangular-lattice description.

\begin{definition}[Inside approximation]\label{def:approximation}
Let $D\subset\R^2$ be a bounded simply connected domain with $C^2$ Jordan
boundary.  For each mesh $\delta$ in a sequence tending to zero, let
$\gamma_\delta$ be a simple cycle in $\delta\hexlat$, let $D_\delta$ be its
bounded interior, and include every lattice vertex and edge on or inside
the cycle.  We require $D_\delta\subset D$, that every compact subset of
$D$ is eventually contained in $D_\delta$, and that there are
parametrizations
\[
 \gamma_\delta:S^1\longrightarrow\partial D_\delta,
 \qquad \gamma:S^1\longrightarrow\partial D,
 \qquad \|\gamma_\delta-\gamma\|_\infty\longrightarrow0.
\]
Here each parametrization is a homeomorphism onto the indicated curve.
Write $F_\delta$ for the hexagonal faces enclosed by $\gamma_\delta$, and
$\partial F_\delta$ for those adjacent to the boundary cycle.
\end{definition}

Fix $x\in[1/\sqrt2,1]$.  An admissible height is a function
$h:F_\delta\to\Z$ such that
\[
 h=0\quad\text{on }\partial F_\delta,
 \qquad |h(u)-h(v)|\le1\quad\text{whenever }u\sim v.
\]
Its probability is
\begin{equation}\label{eq:model}
 \Prob_{\delta,x}(h)=\frac{x^{N_\delta(h)}}{Z_{\delta,x}},
 \qquad
 N_\delta(h)=\sum_{\{u,v\}\subset F_\delta:\,u\sim v}
                 \ind_{\{h(u)\ne h(v)\}}.
\end{equation}
The sum counts each unordered adjacent pair once.  The admissible set is
finite, since the height at any face is bounded in absolute value by its
graph distance to $\partial F_\delta$.  It contains the zero function, so
$Z_{\delta,x}$ is finite and positive.  Height reflection preserves
\eqref{eq:model}; in particular, $\E h_\delta(u)=0$.

Extend $h_\delta$ as a constant on each enclosed face and by zero outside
$D_\delta$.  For $f\in C_c^\infty(D)$ we always use the Euclidean area
pairing
\begin{equation}\label{eq:area-pairing}
 h_\delta(f)=\int_D h_\delta(z)f(z)\dd z
 =\sum_{u\in F_\delta} w_{\delta,f}(u)h_\delta(u),
 \qquad w_{\delta,f}(u)=\int_{\text{face at }u}f(z)\dd z.
\end{equation}
For small enough $\delta$, these weights sum to $\int_D f$ exactly.

Let $\Phi_D$ denote the real zero-Dirichlet Gaussian free field with
covariance
\begin{equation}\label{eq:gff-covariance}
 \E[\Phi_D(f)\Phi_D(g)]
 =\iint_{D\times D} f(z)G_D(z,w)g(w)\dd z\dd w,
 \qquad G_D=(-\Delta_D)^{-1}.
\end{equation}
The Laplacian is the Euclidean Laplacian.  Thus the singular part of
$G_D(z,w)$ is $(2\pi)^{-1}\log(1/|z-w|)$.

\begin{theorem}\label{thm:main}
For every fixed $x\in[1/\sqrt2,1]$ there exists
$\sigma(x)\in(0,\infty)$, independent of $D$ and of the inside approximation
in Definition~\ref{def:approximation}, such that
\[
 \frac{h_\delta}{\sigma(x)}\ \Longrightarrow\ \Phi_D
 \qquad\text{in }\mathcal D'(D).
\]
In particular, the fields are tight as random distributions and, for every
$k\ge1$ and $f_1,\ldots,f_k\in C_c^\infty(D)$, the vector
$(h_\delta(f_i)/\sigma(x))_{i=1}^k$ converges to
$(\Phi_D(f_i))_{i=1}^k$.  All mixed moments of these vectors converge.
\end{theorem}

There is no logarithmic rescaling of the amplitude.  The logarithmic growth
of point variances \cite[Section~1.2.4]{GL2025} is compatible with finite fluctuations after smooth
spatial averaging.  The theorem is for each fixed $x$; it neither specifies
a formula for $\sigma(x)$ nor asserts that this constant is independent of
$x$.

The loop description is exact.  At each interior vertex of $\hexlat$, the
three incident face heights have range at most one.  Consequently either
zero or two incident edges separate unequal heights.  The occupied edges
therefore form mutually vertex-disjoint loops away from the boundary
cycle.  Conversely, such a loop family, together with an independent
choice of one of its two orientations for each loop, determines an
admissible height by summing its signed nesting increments from the zero
exterior.  Thus its un-oriented weight is
\begin{equation}\label{eq:loop-weight}
 2^{\#\text{loops}}x^{\#\text{occupied edges}}.
\end{equation}
Theorem~\ref{thm:main} concerns the field law; no convergence claim for the
individual loops or their ensemble is needed.

\subsection{Context and significance}

Schramm's problem collection asks for Gaussian free field behavior of the
uniform odd-integer Lipschitz height on the triangular lattice with two
boundary arcs, and separately asks about its interface
\cite[Problem~2.2]{Schramm2007}.  That formulation has a different boundary
condition and height convention from \eqref{eq:model}.  It motivates the
field question studied here.  The real-valued model in his Problem~2.3 is
another distinct model.

The loop description belongs to the classical solid-on-solid viewpoint:
Nienhuis relates honeycomb loop weights to integer heights on the
triangular lattice \cite[p.~1062]{Nienhuis1982}.
The predicted critical curve discussed in
\cite[Section~1.4.1]{GL2025} gives $x_c(2)=1/\sqrt2$.
Throughout this paper, the phrase \emph{predicted critical endpoint}
refers to that prediction; a sharp description of the complete phase
diagram is not part of the theorem.

At $x=1/\sqrt2$, macroscopic loops and the resulting logarithmic height
variances were established by Duminil-Copin, Glazman, Peled and Spinka
\cite[Theorem~1 and Section~1.2]{DGPS2021}.
Glazman and Manolescu established logarithmic height variances for the uniform
integer Lipschitz model and developed its two-spin representation and
associated correlation and crossing estimates \cite{GM2021}.
Glazman and Lammers proved delocalization and logarithmic variance
estimates in a broader loop-parameter regime, and related these estimates
to the predicted Gaussian free field limit
\cite[Section~1.2.4]{GL2025}.  More recently, Glazman, Harel and Zelesko
gave another proof of delocalization: for $n=2$ and every
$x\in[1/\sqrt2,1]$, each translation-invariant loop Gibbs measure
has infinitely many loops surrounding every face
\cite[Theorem~2]{GHZ2025}.  These are roughness results, distinct from
the field limit proved here.

The proof develops the reflection and moment method of
the companion article \cite[Sections~5, 7, 9, and 10]{UniformIntegerGFF}
for the weighted family \eqref{eq:model}.  That article treats uniform
odd integer heights with two boundary arcs.  Table~\ref{tab:integer-models}
records the different inputs and conclusions.
The two-spin representation of Glazman--Manolescu and the weighted
percolation augmentation of Glazman--Lammers provide the starting point
for the conditional cuts \cite{GM2021,GL2025}.  The extension requires
order inequalities after restricted observations and estimates for
moments conditioned on rare weighted pinning events.  We prove these
statements below and give the ensuing spectral, boundary, and moment
arguments in full.

\begin{table}[!htbp]
\centering
\small
\renewcommand{\arraystretch}{1.22}
\begin{tabular}{p{0.18\textwidth}p{0.35\textwidth}p{0.35\textwidth}}
\hline
 & Uniform two-arc companion & Present weighted model \\
\hline
Height & Odd integers, adjacent differences $0$ or $2$ &
Integers, adjacent differences at most $1$ \\
Law & Uniform; $H=(h-1)/2$ gives unit-weight increments &
Weight $x$ for each unequal adjacent pair \\
Boundary & Two marked arcs with values $+1$ and $-1$ &
Zero on the boundary layer of an inside approximation \\
Cuts & Uniform spin cuts and loop estimates &
The $x$-dependent open-cell table and complete first-color/status histories \\
Normalization & An absolutely convergent finite-volume series &
An implicit positive constant $\sigma(x)$ for each fixed $x$ \\
Limit & Harmonic boundary profile plus a Dirichlet GFF &
A zero-Dirichlet GFF \\
\hline
\end{tabular}
\caption{The two integer-height articles retain separate model-specific
proofs.  The common reflection, pin-motion, and moment strategy does not
identify their normalization constants or extend the weighted theorem
to two-arc boundary data.}
\label{tab:integer-models}
\end{table}
\FloatBarrier

Reflection positivity provides a classical route from spatial symmetries
to positive transfer operators; see Fr\"ohlich--Israel--Lieb--Simon
\cite{FILS1978} and the scalar spectral discussion of Usui
\cite[Sections~3--4]{Usui2012}.  Here the original edge interaction
factorizes across a row of sites.  We prove the associated positive
normal transfer, including its phase-augmented form, directly.
Appendix~\ref{sec:scalar-spectra} constructs the scalar measures needed
for its Cauchy representation and for analytic translation.

The six-vertex GFF theorem of Duminil-Copin, Kozlowski, Lammers, and
Manolescu provides substantial methodological precedents through transfer
spectra, analytic translation, and harmonic multipoint functions
\cite[Section~4.5; Lemma~7.4; Section~7.2]{DCKLM2026}.  Their result concerns a
full-plane zero-slope state on the square lattice for
$a=b=1$ and $c\in[\sqrt3,2]$.  Its covariance normalization uses
surface tension and Bethe ansatz.  Our plane coefficient is fixed by a
three-mirror angular identity and a signed cutoff limit; the passage to
the actual bounded-domain boundary additionally uses weighted rare-pin
comparisons and boundary attachment.  The final collision-to-moment
argument also has a precedent in Kenyon's identification of dimer-height
correlations \cite[Lemma~3.1 and Proposition~3.2]{Kenyon2001}.  In the present model the
collision coefficient comes from conditional local matching and
logarithmic harmonic expansions.

A related six-vertex result \cite[Theorem~1.1]{BalancedSixVertexGFF2026}
establishes the GFF limit for $a=b=1$ and every fixed $c\in(0,2]$, in the
zero-slope plane law defined by its prescribed iterated limit of balanced tori.

The conclusion strengthens logarithmic variance estimates to a limit
law for the entire distribution-valued field.  Identifying a covariance
is only one part of this conclusion: the proof also establishes Gaussian
higher moments, the zero boundary condition, and tightness.  The parameter
interval includes both the uniform height model $x=1$ and its predicted
critical endpoint $x=1/\sqrt2$.

\subsection{Structure of the proof}

\paragraph{Exact weighted cuts and conditional order.}
Section~\ref{sec:discrete} encodes a height modulo four by two spins and
augments each upward triangular cell with a status.  An open circuit fixes
the first spin while leaving the second spin free.  The resulting cell
factorization gives an exact conditional cut.  The order statement is
proved for histories which record every first-spin site of an observed
cell.  This restriction is essential: conditioning on a status alone need
not preserve positive association.  Stopped explorations then give local
mixing and multiplicative comparisons for separated pins.

\paragraph{Uniform moments and the plane coefficient.}
Section~\ref{sec:moments} proves moment bounds for smeared heights,
including irregular domains exposed by explorations.  A discrete midpoint
inequality of Klartag--Lehec \cite[Theorem~1.4]{KlartagLehec2019}, in the
nonnegative-function form of \cite[Theorem~3]{GRST2019}, amplifies central
probability bounds to all moments.  We derive the finite-support zero
case directly from \cite[Theorem~8]{GRST2019}, then verify the
multidimensional application and rounding conditions explicitly.
In Section~\ref{sec:plane},
site reflection positivity gives a positive normal transfer operator.
Its covariance is a mixture of Cauchy kernels in normal frequency.
An angular identity using all three mirror directions forces every
scaling limit of the spectral measure to be $c\,\dd p/|p|^2$.
A signed cutoff identity proves uniqueness of $c$, and a two-disk variance
bound proves $c>0$.  With our Fourier convention,
\begin{equation}\label{eq:normalization}
 v=(2\pi)^2c,\qquad \sigma(x)=\sqrt v.
\end{equation}

\paragraph{Pins, reflection, and the true boundary.}
Section~\ref{sec:pins} compares the field with upper and lower pinning
brackets.  Uniform normalized moment bounds allow analytic movement of
separated pins in three directions.  Averaging their phases and then
closing the gaps transfers the reflection-null Laplace insertion from
the plane to these brackets.  Section~\ref{sec:boundary} proves the
boundary attachment estimate: a synthetic boundary law gives a favorable
level path, and a deterministic attachment argument connects it to the
required boundary segment.  Iterating in tubes yields boundary arches,
makes the bracket means vanish, and controls boundary reference averages.

\paragraph{Harmonic moments and the full field law.}
Section~\ref{sec:limit} first establishes distributional tightness and
uniform integrability.  The ordered coupling and the vanishing bracket
means then identify the bracket limits with the true field and transfer
the insertion identity to it.  The insertion and boundary references
make every limiting moment distribution separately harmonic away from
collisions.  Logarithmic moment bounds exclude distributions supported
on the diagonals.  Local comparison with the plane fixes the coefficient
of every collision singularity.  Subtracting the corresponding Green
functions gives Wick's recursion, which determines the Gaussian law.
This completes the proof of Theorem~\ref{thm:main}.

\subsection{Conventions}

The parameter $x$ is fixed.  Constants may depend on $x$ and on a stated
fixed macroscopic arrangement; this dependence is not asserted to be
uniform over the parameter interval.  The symbols $C,c,\alpha$ denote
positive constants whose value may change between occurrences.  A test
at radius $r$ has support in a ball of that radius and magnitude
$O(r^{-2})$.  A zero-total test has integral zero.  On the full plane we
use
\[
 \widehat f(p)=\int_{\R^2}e^{-ip\cdot z}f(z)\dd z,
 \qquad
 \langle f,(-\Delta)^{-1}g\rangle
 =\frac1{(2\pi)^2}\int_{\R^2}
       \frac{\overline{\widehat f(p)}\widehat g(p)}{|p|^2}\dd p
\]
for real smooth compactly supported zero-total tests.  The symbol $H$
denotes a generic integer lift in the discrete representations; it is
distinct from the hexagonal lattice $\hexlat$.

Every argument involving collars, windows, disks, or tubes first fixes
their positive macroscopic widths and separations, then sends
$\delta\downarrow0$.  Limits that close gaps, shrink windows, or approach
the boundary are taken afterward.  A bound for fixed geometry is not
used uniformly for a degenerating geometry.  Explorations are stopped
before inspecting the interior region whose conditional Gibbs law is
subsequently used.
